\section{RESPUESTAS}

\subsection{I)}

{\noindent\large\begin{align*}\displaystyle F(x, y, z) &= \sum(2, 3, 4, 5)\end{align*}}
\subsubsection{Obtener la mínima expresión booleana}
\begin{center}
\begin{tabular}{c|ccc|c}
\rowcolor[HTML]{FFCCC9} 
\cellcolor[HTML]{9AFF99}m & x & y & z & \cellcolor[HTML]{FFFFC7}F(x,y,z) \\ \hline
0                         & 0 & 0 & 0 &                                  \\
1                         & 0 & 0 & 1 &                                  \\
2                         & 0 & 1 & 0 & 1                                \\
3                         & 0 & 1 & 1 & 1                                \\
4                         & 1 & 0 & 0 & 1                                \\
5                         & 1 & 0 & 1 & 1                                \\
6                         & 1 & 1 & 0 &                                  \\
7                         & 1 & 1 & 1 &                                 
\end{tabular}
\end{center}
\noindent Los ''1'' de la función pasan, mientras que los ''0'' se los niega. Es una suma de productos de las entradas.
{\large\begin{align*}
    F(x,y,z) &= \overline{x}\cdot y\cdot\overline{z} + \overline{x}\cdot y\cdot z + x\cdot\overline{y}\cdot\overline{z} + x\cdot\overline{y}\cdot z
\end{align*}}
\noindent Simplificación con Karnaugh:
\begin{center}
    \begin{karnaugh-map}[2][4][1][$z$][$x\qquad y$]
    \minterms{2,3,4,5}
    \autoterms[0]
    \implicant{2}{3}\implicant{4}{5}
\end{karnaugh-map}
\end{center}
\begin{itemize}[nosep] \item Primer bloque: ($m_2$ y $m_3$) cambia $z$ y hay que negar $x$.
\item Segundo bloque: ($m_4$ y $m_5$) cambia $z$ y hay que negar $y$.
\item La función simplificada queda:\end{itemize}
    {\large\begin{align*}
    F(x,y,z) &= \overline{x}\cdot y + x\overline{y}
\end{align*}}

\noindent Si se hace una tabla de verdad con las entradas $x$ e $y$ (las funcionales):

\begin{center}
\begin{tabular}{c|cc|c}
\rowcolor[HTML]{FFCCC9} 
\cellcolor[HTML]{9AFF99}m & x & y & \cellcolor[HTML]{FFFC9E}F \\ \hline
0                         & 0 & 0 & 0                         \\
1                         & 0 & 1 & 1                         \\
2                         & 1 & 0 & 1                         \\
3                         & 1 & 1 & 0                        
\end{tabular}


\end{center}

\noindent es identico a una tabla de verdad de $XOR$, por lo que la función se puede simplificar aún más a:

{\large\begin{align*} F(x,y,z) &= x\oplus y\end{align*}}
\subsubsection{Simulación}

\imagen[]{10cm}{./imagenes/ejercicio1Falstad.png}

\subsubsection{Descripción en Verilog}
\begin{lstlisting}
module descripcionVerilogEjercicio1(
    input inputX,
    input inputY,
    input inputZ,
    output outputFunction
    );
	 assign outputFunction = inputX ^ inputY;
endmodule
\end{lstlisting}
\subsubsection{RTL output}

\imagen[Bloque RTL]{15cm}{./imagenes/rtlEjercicio1Verilog.png}
\imagen[Implementación de las compuertas]{15cm}{./imagenes/rtlCompuertasEjercicio1Verilog.png}

\subsubsection{Simulación temporal}


